\section{Conclusion}
\label{sec:conclusion}

Four findings summarize this study. To first order, oxide-clad strip waveguides cannot report their sidewall angle, because their mirror
symmetry makes a tilt equivalent to a width change. One rib from a separate etch leaves the angle unidentifiable. Two
widths of the same rib restore it, and simulations indicate a bound of \qty{0.39}{\degree} for five strip rings and two rib rings on a
\qty{50}{nm} slab at a \qty{1}{pm} wavelength precision, reached by nonlinear fits of data from direct mode solutions when the structures share one set of
process parameters. Fabrication consistency is the practical limit. Random offsets at the levels reported for real
processes raise the bound of one copy of the set to several degrees, and joint fits of $N$ copies recover it as
$1/\sqrt{N}$, with tens of copies needed to return within twice the shared-parameter bound. Offsets that follow a design do not average, and
a fabricated test chip must show whether they stay within a few hundredths of a nanometre.
